\documentclass[10pt,a4paper]{amsart}
\usepackage[margin=19mm]{geometry}
\usepackage{amsmath,amssymb,mathtools}
\usepackage[scr=rsfs,cal=euler]{mathalfa}
\usepackage{xfrac}
\usepackage[unicode,hidelinks,bookmarks=false]{hyperref}
\newcommand{\ie}{i.e.\ }
\newcommand{\cf}{cf.\ }
\newcommand{\resp}{respectively\ }
\newcommand{\Subsection}{\subsection}
\providecommand{\nobreakdash}{\nobreak\hbox{-}}
\setlength{\emergencystretch}{2em}
\multlinegap0pt
\usepackage{bm}
\usepackage{bbm}
\usepackage{dsfont}
\usepackage{xspace}
\usepackage{cancel}
\usepackage{mathtools}
\usepackage{amsthm}
\makeatletter
\@ifundefined{thm}{\newtheorem{thm}{Theorem}}{}
\makeatother
\title[Designing optimal dual frames for $\ell^p-$average error opt]{Designing optimal dual frames for $\ell^p-$average error optimization}
\author{Shankhadeep Mondal, Deguang Han, R. N. Mohapatra}
\date{Source: arXiv:2508.07158v1 (CC BY 4.0)}
\begin{document}
\maketitle

\noindent\emph{Source and licence.} Designing optimal dual frames for $\ell^p-$average error optimization. Version arXiv:2508.07158v1 (CC BY 4.0), distributed under the Creative Commons Attribution 4.0 International licence, as recorded in the arXiv licence metadata for this exact version and shown on the abstract page https://arxiv.org/abs/2508.07158v1. Full rights notice: licenses/arxiv-2508.07158-CC-BY-4.0.txt.\par\smallskip

\noindent\textit{Editorial scope.} Counted unit: the Thm whose source label is \texttt{thm3point4} in the article (source0.tex lines 481--484, source label \texttt{thm3point4}), delivered together with the article's own complete original proof of it (source0.tex lines 487--539, one \texttt{proof} environment of 53 lines). No mathematics is written, repaired or re-derived: statement and proof are the article's own text line for line. Required context retained uncounted: eqn4.6 (lines 295--301); eqn4point8 (lines 464--476). Every definition, display and label of those lines is retained unchanged. Omitted from this package and not redistributed: 1--294, 302--463, 477--480, 485--486, 540--1144. Nothing inside a retained excerpt is omitted or altered: every retained body is delivered exactly as the article's lines are written, so no omission receipt of any kind accompanies this package. Citations: the retained excerpts carry no citation command at all, so no bibliography entry of the article is transcribed and no reference is redistributed. Reference adaptation: none was needed and none was made; every reference inside the delivered unit resolves inside it, so no unresolved reference or citation is produced. Printed slips kept literal: no printed word, symbol, hypothesis or number of the counted statement or of its proof was altered. Layout and class: the article's own class is \texttt{elsarticle} with packages including amssymb, amsmath, amsthm, inputenc, marginnote, amsfonts, textcomp, color, soul, placeins, tikz, makecell, graphicx, hyperref; the wrapper is the lane's pinned \texttt{amsart} editorial wrapper at 10pt on A4, which restates the theorem-like environments the retained text needs and adds the article's own macro definitions that the retained text needs (none), with extra wrapper packages (bm, bbm, dsfont, xspace, cancel, mathtools). The delivered numbering is the unit's own. Assets: no retained span contains a figure environment, a graphics-inclusion command, a PDF-page inclusion, a table or a TikZ picture; no image file is copied into this package, the package declares no asset dependency and no image file is redistributed with it. Licence: version arXiv:2508.07158v1 of this article is distributed under the Creative Commons Attribution 4.0 International licence, as recorded in the arXiv licence metadata for this exact version and shown on the abstract page https://arxiv.org/abs/2508.07158v1. A full rights notice is delivered at licenses/arxiv-2508.07158-CC-BY-4.0.txt. Characters: every retained excerpt is pure ASCII, so the delivered engine, XeLaTeX with its default Latin Modern text font, reports no missing character.
\begin{align}\label{eqn4.6}
\mathrm{AE}_{\mathfrak{F}}^{(1),p}(F, G)
&= \left\{ \frac{1}{N} \sum_{i=1}^N \|f_i\|^p \cdot \|g_i\|^p \right\}^{1/p} \notag \\
&\geq \left( \frac{1}{N} \sum_{i=1}^N \|f_i\| \cdot \|g_i\| \right) \notag \\
&\geq \frac{1}{N} \sum_{i=1}^N \left| \langle f_i, g_i \rangle \right| \notag \\
&\geq \frac{n}{N}.
\end{align}

\begin{align} \label{eqn4point8}
\left\|E_{\Lambda, F, G'}\right\|_{\mathfrak{F}} 
&= \left\|T_{G'}^* D T_F\right\|_{\mathfrak{F}} \nonumber \\
&= \sqrt{\mathrm{tr}\left(D T_F T_F^* D T_{G'} T_{G'}^*\right)} \nonumber \\
&= \sqrt{ \sum_{r=1}^m \langle g'_{i_1}, g'_{i_r} \rangle \langle f_{i_r}, f_{i_1} \rangle 
      + \sum_{r=1}^m \langle g'_{i_2}, g'_{i_r} \rangle \langle f_{i_r}, f_{i_2} \rangle 
      + \cdots 
      + \sum_{r=1}^m \langle g'_{i_m}, g'_{i_r} \rangle \langle f_{i_r}, f_{i_m} \rangle } \nonumber \\
&= \sqrt{ \sum_{r=1}^m \|g'_{i_r}\|^2 \|f_{i_r}\|^2 
      + \sum_{\substack{j,k=1 \\ j \neq k}}^m \langle g'_{i_j}, g'_{i_k} \rangle \langle f_{i_k}, f_{i_j} \rangle } \nonumber \\
&= \sqrt{ \sum_{r=1}^m \|g'_{i_r}\|^2 \|f_{i_r}\|^2 
      + 2\, \mathrm{Re} \left( \sum_{\substack{j,k=1 \\ j > k}}^m \langle g'_{i_j}, g'_{i_k} \rangle \langle f_{i_k}, f_{i_j} \rangle \right) }.
\end{align}

     \begin{thm}\label{thm3point4}
Let \( F = \{f_i\}_{i=1}^N \) be a uniform tight frame for the Hilbert space \( \mathcal{H}_n \), such that \( |\langle f_i, f_j \rangle| \) is constant for all \( i \neq j \). Then, for any \( 1 \leq m \leq N \) and \( p > 0 \), the average Frobenius-norm error over all \( m \)-erasure patterns satisfies $\mathrm{AE}_{\mathfrak{F}}^{(m),p}(F) = \sqrt{m\left(\frac{n}{N}\right)^2 + \frac{m(m-1)(nN - n^2)}{N^2(N-1)}}. $
Moreover, the canonical dual frame \( S_F^{-1} F \) is the unique optimal dual of \( F \) under the Frobenius norm for any \( m-\)erasure. 
\end{thm}

   \begin{proof}
We first prove the result for the case \( m = 1 \). Let \( F \) be a uniform tight frame with frame bound \( A \). Then each frame vector satisfies \( \|f_i\| = \sqrt{\frac{An}{N}} \) for all \( 1 \leq i \leq N \). The canonical dual frame is \( S_F^{-1}F = \left\{ \frac{1}{A}f_i \right\}_{i=1}^N \). Then the average Frobenius-norm error becomes
\[
\mathrm{AE}_{\mathfrak{F}}^{(1),p}(F, S_F^{-1}F) = \left\{ \frac{1}{N} \sum_{i=1}^N \left\|f_i\right\|^p \left\|\frac{1}{A} f_i \right\|^p \right\}^{1/p} = \left\{ \frac{1}{N A^p} \sum_{i=1}^N \|f_i\|^{2p} \right\}^{1/p}.
\]
Since \( \|f_i\| = \sqrt{\frac{An}{N}} \), we obtain
\[
\mathcal{F}_1^p(F, S_F^{-1}F) = \left\{ \frac{1}{N A^p} \cdot N \left( \frac{An}{N} \right)^p \right\}^{1/p} = \frac{n}{N}.
\]
Hence by \eqref{eqn4.6}, the canonical dual $S_F^{-1}F$ is a $1-$erasure optimal dual of \( F \) under the Frobenius norm.

Now consider the case \( m > 1 \). Let \( G \in \zeta_{\mathfrak{F}}^{(m),p}(F) \). Then \( G \in \zeta_{\mathfrak{F}}^{(1),p}(F) \), and thus,
\[
\mathrm{AE}_{\mathfrak{F}}^{(1),p}(F, G)=\left\{ \frac{1}{N} \sum_{i=1}^N \|f_i\|^p \|g_i\|^p \right\}^{1/p} = \frac{n}{N}.
\]
Now,
\[
N \left(\frac{n}{N}\right)^p = \sum_{i=1}^N \|f_i\|^p \|g_i\|^p \geq N \left( \frac{\sum_{i=1}^N \|f_i\| \|g_i\|}{N} \right)^p \geq N \left( \frac{\sum_{i=1}^N |\langle f_i, g_i \rangle|}{N} \right)^p \geq N \left( \frac{n}{N} \right)^p.
\]
Thus, equality holds throughout, which leads to $\|f_i\| \|g_i\| = |\langle f_i, g_i \rangle| = \frac{n}{N}, \quad \forall i.$ For a frame $F=\{f_i\}_{i=1}^N$ and a dual $G=\{g_i\}_{i=1}^N,$ we have

\[
\sum_{i,j=1}^N \langle g_i, g_j \rangle \langle f_j, f_i \rangle = \sum_{i=1}^N \left\langle g_i, \sum_{j=1}^N \langle f_i, f_j \rangle g_j \right\rangle = \sum_{i=1}^N \langle g_i, f_i \rangle = n.
\]
Also, since \( \|f_i\| \|g_i\| = \frac{n}{N} \), we derive:
\begin{align} \label{equation3point8}
2\,\text{Re} \left( \sum_{\substack{i,j=1 \\ i > j}}^N \langle g_i, g_j \rangle \langle f_j, f_i \rangle \right) = \sum_{i \neq j} \langle g_i, g_j \rangle \langle f_j, f_i \rangle = n - \frac{n^2}{N}.
\end{align}

Now, from equation~\eqref{eqn4point8}, 
\begin{align*}
\mathrm{AE}_{\mathfrak{F}}^{(m),p}(F, G)
&= \left\{ \frac{1}{\binom{N}{m}} \sum_{D \in \mathcal{D}^{(m)}} \left( \sum_{r=1}^m \|g_{i_r}\|^2 \|f_{i_r}\|^2 + \sum_{\substack{j,k=1 \\ j \neq k}}^m \langle g_{i_j}, g_{i_k} \rangle \langle f_{i_k}, f_{i_j} \rangle \right)^{p/2} \right\}^{1/p} \\
&= \left\{ \frac{1}{\binom{N}{m}} \sum_{D \in \mathcal{D}^{(m)}} \left( m \left( \frac{n}{N} \right)^2 + \sum_{\substack{j,k=1 \\ j \neq k}}^m \langle g_{i_j}, g_{i_k} \rangle \langle f_{i_k}, f_{i_j} \rangle \right)^{p/2} \right\}^{1/p} \\
&\geq \left( m \left( \frac{n}{N} \right)^2 + \frac{m(m-1)}{N(N-1)} \sum_{i \neq j} \langle g_i, g_j \rangle \langle f_j, f_i \rangle \right)^{1/2} \\
&= \sqrt{ m\left( \frac{n}{N} \right)^2 + \frac{m(m-1)(nN - n^2)}{N^2(N-1)} }.
\end{align*}

Thus, $\mathrm{AE}_{\mathfrak{F}}^{(1),p}(F) \geq \sqrt{ m\left( \dfrac{n}{N} \right)^2 + \dfrac{m(m-1)(nN - n^2)}{N^2(N-1)} }.$ Also,
\begin{align*}
\mathrm{AE}_{\mathfrak{F}}^{(1),p}(F, S_{F}^{-1}F)
&= \left\{ \frac{1}{\binom{N}{m}} \sum_{D \in \mathcal{D}^{(m)}} \left( \sum_{r=1}^m \frac{1}{A^2} \|f_{i_r}\|^4 + \frac{1}{A^2} \sum_{\substack{j,k=1 \\ j \neq k}}^m |\langle f_{i_j}, f_{i_k} \rangle|^2 \right)^{p/2} \right\}^{1/p} \\
&= \sqrt{ m\left( \dfrac{n}{N} \right)^2 + \dfrac{m(m-1)(nN - n^2)}{N^2(N-1)} }.
\end{align*}

Hence, $ \mathrm{AE}_{\mathfrak{F}}^{(m),p}(F) = \sqrt{ m\left( \dfrac{n}{N} \right)^2 + \dfrac{m(m-1)(nN - n^2)}{N^2(N-1)} }.$  It remains to prove the uniqueness. Suppose there exists another optimal dual \( \tilde{G} = \left\{ \frac{1}{A} f_i + u_i \right\}_{i=1}^N \) satisfying \(  \mathrm{AE}_{\mathfrak{F}}^{(1),p}(F, \tilde{G}) = \mathrm{AE}_{\mathfrak{F}}^{(1),p}(F) = \frac{n}{N} \). Therefore, $\frac{n}{N}=\left\{\frac{1}{N} \sum \limits_{i=1}^N\|f_i\|^p\,\| \tilde{g}_i \| ^p\right\}^{\frac{1}{p}} \geq \dfrac{\sum \limits_{i=1}^N\|f_i\|^p\,\| \tilde{g}_i \|}{N} \geq \frac{n}{N}. $ This leads to $\left\{\frac{1}{N} \sum \limits_{i=1}^N\|f_i\|^p\,\| \tilde{g}_i \| ^p\right\}^{\frac{1}{p}} = \dfrac{\sum \limits_{i=1}^N\|f_i\|\,\| g_i \|}{N} = \frac{n}{N}.$ This is possible if and only if  each $\|f_i\|\,\| \tilde{g}_i \|$ are equal and hence, $\|f_i\|\,\| \tilde{g}_i \| = \|f_i\|\left\| \frac{1}{A}f_i + u_i \right\|= \frac{n}{N},$ for all $1 \leq i \leq N.$ This leads to, $\left\| \frac{1}{A}f_i + u_i \right\| = \sqrt{\frac{n}{AN}},\;1 \leq i \leq N.$ Explicitly, $\frac{1}{A^2}\|f_i\|^2 + \frac{2}{A} Re\langle f_i, u_i \rangle + \|u_i\|^2 = \frac{n}{AN},\;1 \leq i \leq N.$ Thus, $\frac{1}{A^2}\sum\limits_{1 \leq i \leq N} \|f_i\|^2 + \frac{2}{A} \sum\limits_{1 \leq i \leq N} Re\langle f_i, u_i \rangle + \sum\limits_{1 \leq i \leq N} \|u_i\|^2  = \frac{n}{A}.$ Consequently,
 \begin{align}\label{equation9}
 	\frac{2}{A} \sum\limits_{1 \leq i \leq N} Re\langle f_i, u_i \rangle + \sum\limits_{1 \leq i \leq N} \|u_i\|^2 = 0.
 \end{align}
As $\tilde{G}$ is a dual of $F,$  $\sum\limits_{1\leq i \leq N} \langle f,f_i \rangle u_i =0,$ for all $f \in \mathcal{H}_n.$ This implies, $T_{U}^* T_F = 0,$ where $U= \{u_i\}_{i=1}^N.$ Therefore, $0 = tr(T_{U}^* T_F) = tr(T_F T_{U}^*) = \sum\limits_{1\leq i \leq N} \langle f_i,u_i \rangle. $ So, $Re\left( \sum\limits_{1\leq i \leq N} \langle f_i,u_i \rangle \right) =0.$ Hence, by \eqref{equation9}, $\sum\limits_{1\leq i \leq N} \|u_i\|^2 = 0,$ which implies, $u_i = 0,$ for all $1 \leq i \leq N.$ Thus, $\tilde{G} = S_{F}^{-1}F $.


\end{proof}

\end{document}
